\chapter{Szkeleton tervezése}

\section{Fontos változások a korábbi tervekhez képest}
\par{A korábbi elképzeléseinkkel ellentétben a komplex Csomópont objektumok szerepét egyszerű betűkkel jelölt azonosítók veszik át, amelyeket az Út osztály tárol el attribútumként, meghatározva az adott szakasz kezdő- és végpontját. Erre a változtatásra azért volt szükség, mert a csomópontnak önmagában nem volt elegendő önálló felelőssége a rendszerben. Hosszas megbeszélések után arra jutottunk, hogy a korábbi, túl összetett felépítés a jövőbeni implementáció során súlyos komplikációkhoz, illetve a tantárgy szabályaival való ütközéshez vezetett volna, így az azonosító alapú hivatkozás stabilabb alapot biztosít az úthálózat gráfjának kezeléséhez.}

\section{A szkeleton modell valóságos use-case-ei}
\subsection{Use-case diagram}
\diagram{docs/img/uml_use_case/uc1}{Biokerozin töltése}{16cm}
\diagram{docs/img/uml_use_case/uc2}{Biokerozin töltése}{12cm}
\clearpage

\subsection{Use-case leírások}

\begin{use-case}
    {Biokerozin töltés}
    {Biokerozin töltés a Sárkány fejbe a TakarítóManager által kezdeményezve}
    {TakarítóManager}
    Hasonlóan a sóhoz, a sárkány fej működtetéséhez szükséges biokerozin is a Telephelyen vételezhető. A takarító szerepet játszó játékos pénzért cserébe tölti fel a munkagép üzemanyagtartályát, biztosítva a gázturbina további működését. A rendszer a tranzakció sikere után frissíti a hókotró biokerozin-készletét, lehetővé téve a takarítás végrehajtását.
\end{use-case}

\begin{use-case}
    {Busz haladása és fizetés}
    {Busz haladása és fizetés a BuszsoforManager által kezdeményezve}
    {BuszsoforManager}
    A BuszsoforManager szerepkörben tevékenykedő játékos kezdeményezi a busz léptetését az úthálózaton. A jármű a megadott cél felé halad, miközben a rendszer ellenőrzi, hogy az adott kereszteződés egy Végállomás-e. Amennyiben a busz végállomásra ér, a program regisztrálja a megtett kört, növeli a teljesített utak számát, és a játékos vagyona gyarapodik a menetdíjból származó bevétellel.
\end{use-case}

\begin{use-case}
    {Hókotró fejcsere}
    {Hókotró fejének cseréje a TakarítóManager által kezdeményezve}
    {TakarítóManager}
    A TakarítoManager a Telephelyen kezdeményezi a Hókotróra szerelt munkaeszköz cseréjét. A folyamat során a rendszer ellenőrzi, hogy a Jármű a telephelyen tartózkodik-e, majd a régi fejet leszerelik és a raktárba helyezik. Ezt követően, ha a Játékos rendelkezik a felszerelni kívánt új fejjel, az rákerül a gépre, és törlődik a raktárkészletből, hiszen onnantól az lesz az aktív fej a hókotrón.
\end{use-case}

\begin{use-case}
    {Fej vásárlása}
    {Új hókotró fej vásárlása a TakarítóManager által kezdeményezve}
    {TakarítóManager}
    A TakarítoManager a Telephelyen speciális takarítóeszközöket (például Sárkány- vagy Sószóró fejet) vásárolhat a hatékonyabb munkavégzés érdekében. A vásárlás során a rendszer ellenőrzi a pénzügyi keretet, majd a vételár kifizetése után az új eszköz bekerül a játékos raktárába. A megvásárolt fejeket a játékos később bármelyik Hókotrójára felszerelheti a telephelyen.
\end{use-case}

\begin{use-case}
    {Hányó fej működése}
    {Hányó fej működése a jármű mozgatása közben}
    {JátékosManager}
    A fej először ellenőrzi, hogy található-e hóréteg vagy feltört jég az úton. Amennyiben van eltávolítandó csapadék, a fej azt nagy erővel az út mellé szórja, így a hó véglegesen kikerül a rendszerből, és nem kerül át szomszédos sávokba. A tömör jégpáncél viszont változatlanul marad a sávon. 
\end{use-case}

\begin{use-case}
    {Havazás szimulációja}
    {Havazás szimulációja minden pillanatban}
    {Tesztelő}
    A környezeti hatások modellezéséért felelős Úthálózat minden körben végrehajtja a havazás folyamatát. A rendszer egyenletesen növeli a hóvastagságot a város minden útszakaszán, kivéve az alagutakat, amelyeket a kialakításuk megvéd a csapadéktól. 
\end{use-case}

\begin{use-case}
    {Játék vége 1 – Elakadt tömegközlekedés}
    {Játék vége 1 – Elakadt tömegközlekedés a városban}
    {Tesztelő}
    A rendszer minden lépés után ellenőrzi a forgalom állapotát és a játék folytathatóságát. Ez a use case akkor következik be, ha a rendszer megállapítja, hogy a városban egyetlen busz sem képes a mozgásra (például balesetek vagy hótorlaszok miatt). Mivel a tömegközlekedés megbénulása a város működésképtelenségét jelenti, a játék ebben a pillanatban vereséggel véget ér.
\end{use-case}

\begin{use-case}
    {Játék vége 2 – Kijárási tilalom}
    {Játék vége 2 – Kijárási tilalom a városban}
    {Tesztelő}
    A rendszer szintén folyamatosan monitorozza a városban bekövetkezett balesetek és úttorlaszok számát. Amennyiben az összecsúszott és mozgásképtelenné vált járművek száma elér egy előre meghatározott kritikus határértéket, “kijárási tilalmat hirdetnek”. Ez a vészhelyzet a szimuláció azonnali leállítását és a játék végét vonja maga után.
\end{use-case}

\begin{use-case}
    {Autó áthaladása és jéggé tömörülés}
    {Autó áthaladása és jéggé tömörülés a sávon}
    {Tesztelő}
    Amikor egy NPC autó vagy egy busz áthalad egy havas útsávon, súlyával tömöríti az ott lévő havat. A sáv minden áthaladáskor növeli a számlálóját, és amint eléri a kritikus (5 jármű) határt, a sávon lévő teljes hómennyiség azonnal jégpáncéllá tömörödik. Ez a folyamat dinamikusan változtatja az útviszonyokat, növelve a balesetveszélyt a forgalom számára.
\end{use-case}

\begin{use-case}
    {Jégtörő fej működése}
    {Jégtörő fej működése a sávon lévő jégpáncél feltörésére}
    {JátékosManager}
    A fej feladata a sávon lévő masszív jégpáncél fizikai feltörése, amelyet “feltört jéggé” (a modellben hóvá) alakít át, de az útról nem távolítja el, így az visszatömörödhet jéggé. 
\end{use-case}

\begin{use-case}
    {Leszórt só}
    {Leszórt só hatása a sáv állapotára}
    {Tesztelő}
    A sószóró fej által az útra juttatott só aktív hatást fejt ki a sáv állapotára. Felolvasztja a havat és a jeget, illetve meggátolja azok újboli kialakulását egy bizonyos ideig. 
\end{use-case}

\begin{use-case}
    {Jármű megcsúszása}
    {Jármű megcsúszása a jeges sávon}
    {Tesztelő}
    Amikor egy jármű (autó vagy busz) jeges sávra lép, a rendszer baleseti kalkulációt végez a globális nehézségi szint és a jégvastagság alapján. Sikertelen kimenetel esetén a jármű megcsúszik, elszenvedi a balesetet és meghatározott ideig mozgásképtelenné válik. Az ilyen jármű eltorlaszolja az adott sávot, akadályozva a forgalom többi résztvevőjét, amíg a mozgásképtelensége le nem jár.
\end{use-case}

\begin{use-case}
    {Sárkány fej működése}
    {Sárkány fej működése a sávon lévő hó és jég eltávolítására}
    {JátékosManager}
    A fej működésének feltétele, hogy a Sárkányfej bekapcsolt állapotban legyen és rendelkezzen elegendő biokerozinnal. Aktív állapotban a gép biokerozint fogyaszt, cserébe az adott sávon minden havat és jégpáncélt azonnal és véglegesen elolvaszt, megtisztítva az utat a forgalom számára.
\end{use-case}

\begin{use-case}
    {Söprő fej működése}
    {Söprő fej működése a sávon lévő hó eltávolítására}
    {JátékosManager}
    A fej lekérdezi az aktuális sáv hóvastagságát, majd a havat és a feltört jeget a menetirány szerinti jobb oldali szomszédos sávba továbbítja. Ha a gép a legszélső sávban halad, a csapadék az út mellé kerül, ahol már nem akadályozza tovább a közlekedést.
\end{use-case}

\begin{use-case}
    {Sószóró fej működése}
    {Sószóró fej működése a sávon lévő hó és jég olvasztására és újraképződésének megakadályozására}
    {JátékosManager}
    Ez a fej működés közben csökkenti a gép sókészletét, és a sávhoz egy speciális só-objektumot rendel. Ez az anyag nemcsak felolvasztja a meglévő havat és jeget, hanem pár körön keresztül megakadályozza az újabb csapadék megmaradását is az adott sávon. 
\end{use-case}

\begin{use-case}
    {Só töltése a sószóró fejbe}
    {Só töltése a sószóró fejbe a TakarítóManager által kezdeményezve}
    {TakarítóManager}
    A hókotrókat vezető játékos a Telephelyen kezdeményezi a hókotró sószóró fejéhez tartozó sókészlet feltöltését. A folyamat során a rendszer ellenőrzi a játékos anyagi fedezetét, majd a megfelelő összeg levonása után a jármű sókészlete maximális szintre áll vissza. Ez a művelet elengedhetetlen a sószóró fej folyamatos és hatékony üzemeltetéséhez a jeges útszakaszokon.
\end{use-case}

\begin{use-case}
    {Új hókotró vásárlása}
    {Új hókotró vásárlása a TakarítóManager által kezdeményezve}
    {TakarítóManager}
    A Játékos a Telephelyen új Hókotrót vásárolhat a vagyona terhére, hogy bővítse a flottáját. A rendszer ellenőrzi a tartózkodási helyet és az anyagi fedezetet. Ha minden feltétel teljesül, levonja a vételárat a játékos pénzéből. Sikeres tranzakció után létrejön az új hókotró objektum, amely azonnal munkára fogható.
\end{use-case}

\section{A szkeleton kezelői felületének terve, dialógusok}
\par{A szkeleton program egy konzolos alkalmazásként fog futni. A felhasználó a standard bemeneten keresztül irányítja a teszteseteket, a kimenet pedig a standard kimeneten jelenik meg. A kiírásokért, a beolvasásokért és az bekezdések kezeléséért egy statikus Szkeleton osztály felel.}
\par{A program elindításakor a képernyőn megjelenik egy főmenü, amely kilistázza a megvalósított use-case-eket. Ahány szekvenciadiagramot tesztelni kell, annyi menüpont lesz elérhető, kiegészítve egy kilépési opcióval. A felhasználó a sorszám begépelésével választhatja ki a futtatni kívánt esetet.}
\par{Példa a menüre és a bemenetre:}
\begin{verbatim}
    Szkeleton program menü:
    1. Use-case 1
    2. Use-case 2
    3. Use-case 3
    .
    .
    .
    0. Kilépés
    Válassz egy opciót: 
\end{verbatim}
\par{A kiválasztott tesztesetbe lépés után a program elkezdi kiírni a metódushívásokat. Ahogy az objektumok egyre mélyebben hívják meg egymás függvényeit, a kiírt szöveget tabulátorokkal egyre beljebb toljuk, hogy látható legyen a hívások hierarchiája.}
\par{A hívásokat jobbra mutató nyíl jelöli, amit a hívott objektum azonosítója, típusa és a metódus neve, paraméterei követnek. A visszatéréseket balra mutató nyíl jelöli, amit a visszatérési típus vagy érték követ.}
\begin{verbatim}
-> h1: Hokotro.takaritSavot(s_akt)
    -> f1: SoproFej.getHo()
    <- mennyiseg: int
    -> s_jobb: Sav.hoHozzaadas(mennyiseg)
    <- void
    -> s_akt: Sav.takaritHatas(sz, h)
    <- void
<- void
\end{verbatim}
\par{Mivel a szkeleton fázisban az osztályok nem tárolnak belső állapotot, az elágazásoknál és a döntési pontokon a program a felhasználótól kérdez, hogy meghatározza a teszt kimenetelét. Ezeket a kérdéseket a Szkeleton osztály teszi fel, jól látható módon elkülönítve a normál kiírástól.}
\par{Példa a döntési pontokra:}
\begin{verbatim}
Elérte az áthaladt autók száma az 5-öt? (I/N):

Elegendo a játékos vagyona (vagyon >= HOKOTRO_ARA)? (I/N):
\end{verbatim}

\section{Szekvencia diagramok a belső működésre}
\diagram{docs/img/uml_sequence/biokerozin}{Biokerozin töltése}{15cm}
\diagram{docs/img/uml_sequence/buszlep}{Busz haladása}{13.5cm}
\diagram{docs/img/uml_sequence/fejcsere}{Hókotró fejcsere}{9cm}
\diagram{docs/img/uml_sequence/fejvasarlas}{Fej vásárlása}{17cm}
\diagram{docs/img/uml_sequence/hanyofej}{Hányó fej}{17cm}
\diagram{docs/img/uml_sequence/havazas}{Havazás}{15cm}
\diagram{docs/img/uml_sequence/jatekvege_elakadt}{Játék vége - elakadt buszok}{17cm}
\diagram{docs/img/uml_sequence/jatekvege_tilalom}{Játék vége - kijárási tilalom}{17cm}
\diagram{docs/img/uml_sequence/jeggetomorodes}{Autó lép, jéggé tömörödés}{12.5cm}
\diagram{docs/img/uml_sequence/jegtorofej}{Jégtörő fej}{17cm}
\diagram{docs/img/uml_sequence/leszortso}{Leszort só hatása}{12cm}
\diagram{docs/img/uml_sequence/osszecsuszas}{Járművek összecsúszása}{11.5cm}
\diagram{docs/img/uml_sequence/sarkanyfej}{Sárkány fej}{15cm}
\diagram{docs/img/uml_sequence/soprofej}{Söprő fej}{17cm}
\diagram{docs/img/uml_sequence/soszorofej}{Sószóró fej}{17cm}
\diagram{docs/img/uml_sequence/sotoltes}{Só töltés}{14cm}
\diagram{docs/img/uml_sequence/ujhokotro}{Új Hókotró}{14.5cm}
\clearpage

\section{Kommunikációs diagramok}
\diagram{docs/img/uml_communication/biokerozin}{Biokerozin töltése}{16cm}
\diagram{docs/img/uml_communication/buszlep}{Busz haladása}{18cm}
\diagram{docs/img/uml_communication/fejcsere}{Hókotró fejcsere}{17cm}
\diagram{docs/img/uml_communication/fejvasarlas}{Fej vásárlása}{17cm}
\diagram{docs/img/uml_communication/hanyofej}{Hányó fej}{16cm}
\diagram{docs/img/uml_communication/havazas}{Havazás}{16cm}
\diagram{docs/img/uml_communication/jatekvege_elakadt}{Játék vége - elakadt buszok}{16cm}
\diagram{docs/img/uml_communication/jatekvege_tilalom}{Játék vége - kijárási tilalom}{16cm}
\diagram{docs/img/uml_communication/jeggetomorodes}{Autó lép, jéggé tömörödés}{16cm}
\diagram{docs/img/uml_communication/jegtorofej}{Jégtörő fej}{16cm}
\diagram{docs/img/uml_communication/leszortso}{Leszort só hatása}{17cm}
\diagram{docs/img/uml_communication/osszecsuszas}{Járművek összecsúszása}{17cm}
\diagram{docs/img/uml_communication/sarkanyfej}{Sárkány fej}{17cm}
\diagram{docs/img/uml_communication/soprofej}{Söprő fej}{16cm}
\diagram{docs/img/uml_communication/soszorofej}{Sószóró fej}{21cm}
\diagram{docs/img/uml_communication/sotoltes}{Só töltés}{16cm}
\diagram{docs/img/uml_communication/ujhokotro}{Új Hókotró}{16cm}
\clearpage